% Einheit 2 von 3 – Einheit 2 (bank/normalverteilung-und-sigma-regeln/e2.jsonl, zusammenbau v0.1)
%% TODO Prüfungswort Einheit 2: nicht in der Marken-Zeile
%% TODO Zeitmarke Einheit 2: Marken-Zeile nicht lesbar
%% TODO Zweigzeile Teil 1: was hier gelernt wird (Satz in Schülersprache aus dem Einheitstitel) – Einheit 2
%% TODO Einheitentitel 2 nicht in der Mappe lesbar
\einheitenkopf[e2]{Einheit 2 von 3 · Einheit 2}
\setcounter{aufgabe}{9}

%% TODO Titel: Ich-kann-Satz fehlt in der Bank (Platzhalter: Kettenname) – Nr. 10
%% TODO Anweisung: ein Satz über der ersten Teilaufgabe fehlt in der Bank – Nr. 10
\begin{aufgabe}{Wahrscheinlichkeiten berechnen}
%% TODO \rechenplatz{n}: Schrittzahl des Grundfalls fehlt in der Bank (nur bei fester Schreibform, 2.3 b)
\begin{teile}
\teil Gesucht ist $P(480 \le X \le 520)$ für die Füllmenge $X$. Wörtlich oder übersetzt? \\ \kreuz{Das Intervall steht wörtlich da.} \\ \kreuz{Das Intervall muss erst übersetzt werden.}
\teil Die Füllmenge $X$ einer Saftflasche in ml ist normalverteilt mit $\mu = 500$ und $\sigma = 4$. $P(496 \le X \le 506)$? \hfill (Abitur 2023 LK) \\ \leerfeld
\teil Die Abbildung zeigt die Dichte der normalverteilten Zufallsgröße $A$ mit dem Maximum bei $30$; es gilt $P(26 \le A \le 34) \approx 68\,\%$. $P(A > 34)$ ohne Rechner? \hfill (Abitur 2022 LK) \\

\normalverteilung{30}{4}
\leerfeld \%
\teil Die Zahl $T$ der Tore einer Liga in einer Saison wird durch die normalverteilte Zufallsgröße $X$ mit $\mu = 52$ und $\sigma = 6$ beschrieben; eine Anzahl steht im Modell für den Bereich von einer halben Einheit darunter bis einer halben Einheit darüber. $P(T = 50)$? \hfill (Abitur 2026 LK) \\ \leerfeld
\teil Ein Laden gibt je volle $4$ € Einkaufswert einen Punkt; die Zahl $C$ der Punkte bei einem Einkauf von $60$ bis $63{,}99$ € wird durch die normalverteilte Zufallsgröße $X$ mit $\mu = 15$ und $\sigma = 2{,}5$ beschrieben. Annahme: Mit mindestens $60\,\%$ weicht $C$ um höchstens $2$ vom Sollwert ab. Beurteile. \hfill (Abitur 2026 LK)
\teil Die Abbildung zeigt die Dichte von $X$ mit dem Erwartungswert $50$. Jemand rechnet: $P(47 \le X \le 50) \approx 3 \cdot 0{,}1 = 0{,}3$, somit $P(|X - 50| > 3) \approx 1 - 2 \cdot 0{,}3 = 0{,}4$. Erläutere die Überlegungen. \hfill (Abitur 2024 LK)

\normalverteilung{50}{4}
\end{teile}
\end{aufgabe}

%% TODO Titel: Ich-kann-Satz fehlt in der Bank (Platzhalter: Kettenname) – Nr. 11
%% TODO Anweisung: ein Satz über der ersten Teilaufgabe fehlt in der Bank – Nr. 11
\begin{aufgabe}{Wahrscheinlichkeiten berechnen – Fehler finden, Begründen}
\begin{teile}
\teil Ich finde den Fehler: $X$ ist normalverteilt mit $\mu = 80$ und $\sigma = 5$. Ole: „Außerhalb von $70$ bis $90$ liegen etwa $4{,}6\,\%$, also ist $P(X > 90) \approx 4{,}6\,\%$.“
\teil Begründe: Wird eine Anzahl durch eine stetige Zufallsgröße beschrieben, braucht man an den Grenzen halbe Schritte.
\end{teile}
\end{aufgabe}
